\section{AI for Science：LLM 引导的科学发现}

\subsection{从符号回归到 LLM 引导进化}

\begin{knowledgebox}{符号回归问题}
\textbf{符号回归}（Symbolic Regression）是根据经验观测数据发现方程的任务。例如，给定火星运动的测量数据，发现开普勒第三定律。

传统方法（如遗传编程 PySR）维护候选表达式种群，通过变异和交叉进化。LLM 可通过提供先验科学知识来增强这一过程。
\end{knowledgebox}

\subsection{LaSR：概念引导的符号回归}

\begin{importantbox}{LaSR（Language-guided Symbolic Regression）框架}
LaSR 联合推断程序假设和高级概念：

$$P(\pi, C \mid D) \propto P(D \mid \pi) \cdot P(\pi \mid C) \cdot P(C)$$

其中：
\begin{itemize}
    \item $P(D \mid \pi)$：程序 $\pi$ 对数据 $D$ 的似然（通过执行程序计算）
    \item $P(\pi \mid C)$：给定概念 $C$ 下程序的概率（LLM 的背景知识）
    \item $P(C)$：概念的先验分布（LLM 的科学知识）
\end{itemize}
\end{importantbox}

LaSR 的工作流程：
\begin{enumerate}
    \item 维护候选程序（方程）种群
    \item \textbf{假设进化}：经典符号变异 + LLM 引导变异（基于概念库）
    \item \textbf{适应度评估}：程序对数据的拟合程度
    \item \textbf{概念抽象}：LLM 将程序池抽象为自然语言高级描述（如"指数增长/衰减"、"幂律"）
    \item \textbf{概念进化}：LLM 组合现有概念生成新概念（如"指数增长" + "温度依赖" $\to$ "Boltzmann 分布"）
    \item 概念库反馈到假设进化，形成闭环
\end{enumerate}

\subsection{实验结果}

\begin{itemize}
    \item 在费曼方程发现任务上，LaSR 优于纯遗传编程和无概念学习的版本
    \item 即使使用本地语言模型（非前沿模型），也能超过纯遗传编程
    \item 用户提示（hints）可提供额外收益
    \item LaSR 发现的库仑定律表达式仅需4步化简即可得到标准形式，而 PySR 得到的等价表达式需要大量化简
\end{itemize}

\subsection{新发现：LLM 缩放定律}

为验证方法不只是"记忆已知公式"，团队用 LaSR 发现了新的 LLM 缩放定律，纳入了 \textbf{few-shot 样本数}参数：

\begin{warningbox}{有趣发现}
大量 few-shot 样本对低能力模型反而有害，但一旦模型能力超过阈值，更多 shots 带来更好表现。将此发现与 Chinchilla 定律结合，得到了更准确的缩放定律。
\end{warningbox}

\subsection{本章小结}

\begin{itemize}
    \item LLM 引导的进化是经验科学的强大工具
    \item LLM 的抽象能力（概念发现）可进一步增强发现过程
    \item 可扩展到视觉高维输入（使用 VLM）
    \item 开放挑战：假设验证、超越自然语言的概念表示、更大搜索空间的扩展
\end{itemize}

